\documentclass[10pt,a4paper]{amsart}
\usepackage[margin=19mm]{geometry}
\usepackage{amsmath,amssymb,mathtools}
\usepackage[scr=rsfs,cal=euler]{mathalfa}
\usepackage{xfrac}
\usepackage[unicode,hidelinks,bookmarks=false]{hyperref}
\newcommand{\ie}{i.e.\ }
\newcommand{\cf}{cf.\ }
\newcommand{\resp}{respectively\ }
\newcommand{\Subsection}{\subsection}
\providecommand{\nobreakdash}{\nobreak\hbox{-}}
\setlength{\emergencystretch}{2em}
\multlinegap0pt
\usepackage{bm}
\usepackage{bbm}
\usepackage{dsfont}
\usepackage{xspace}
\usepackage{cancel}
\usepackage{mathtools}
\usepackage{amsthm}
\makeatletter
\@ifundefined{lemma}{\newtheorem{lemma}{Lemma}}{}
\makeatother
\title[On Sets of Monochromatic Objects in Bicolored Point Sets]{On Sets of Monochromatic Objects in Bicolored Point Sets}
\author{Sujoy Bhore, Konrad Swanepoel}
\date{Source: arXiv:2602.17637v1 (CC BY 4.0)}
\begin{document}
\maketitle

\noindent\emph{Source and licence.} On Sets of Monochromatic Objects in Bicolored Point Sets. Version arXiv:2602.17637v1 (CC BY 4.0), distributed under the Creative Commons Attribution 4.0 International licence, as recorded in the arXiv licence metadata for this exact version and shown on the abstract page https://arxiv.org/abs/2602.17637v1. Full rights notice: licenses/arxiv-2602.17637-CC-BY-4.0.txt.\par\smallskip

\noindent\textit{Editorial scope.} Counted unit: the Lemma whose source label is \texttt{lemma:3} in the article (source0.tex lines 619--622, source label \texttt{lemma:3}), delivered together with the article's own complete original proof of it (source0.tex lines 623--673, one \texttt{proof} environment of 51 lines). No mathematics is written, repaired or re-derived: statement and proof are the article's own text line for line. Required context retained uncounted: none. Every definition, display and label of those lines is retained unchanged. Omitted from this package and not redistributed: 1--618, 674--1025. Nothing inside a retained excerpt is omitted or altered: every retained body is delivered exactly as the article's lines are written, so no omission receipt of any kind accompanies this package. Citations: the retained excerpts carry no citation command at all, so no bibliography entry of the article is transcribed and no reference is redistributed. Reference adaptation: none was needed and none was made; every reference inside the delivered unit resolves inside it, so no unresolved reference or citation is produced. Printed slips kept literal: no printed word, symbol, hypothesis or number of the counted statement or of its proof was altered. Layout and class: the article's own class is \texttt{article} with packages including fontenc, lmodern, amsmath, amsfonts, bm, amsthm, url, float, tikz, comment, smacros; the wrapper is the lane's pinned \texttt{amsart} editorial wrapper at 10pt on A4, which restates the theorem-like environments the retained text needs and adds the article's own macro definitions that the retained text needs (none), with extra wrapper packages (bm, bbm, dsfont, xspace, cancel, mathtools). The delivered numbering is the unit's own. Assets: no retained span contains a figure environment, a graphics-inclusion command, a PDF-page inclusion, a table or a TikZ picture; no image file is copied into this package, the package declares no asset dependency and no image file is redistributed with it. Licence: version arXiv:2602.17637v1 of this article is distributed under the Creative Commons Attribution 4.0 International licence, as recorded in the arXiv licence metadata for this exact version and shown on the abstract page https://arxiv.org/abs/2602.17637v1. A full rights notice is delivered at licenses/arxiv-2602.17637-CC-BY-4.0.txt. Characters: every retained excerpt is pure ASCII, so the delivered engine, XeLaTeX with its default Latin Modern text font, reports no missing character.
\begin{lemma}\label{lemma:3}
Let $n\geq 6$ blue points on a circle, and $n$ red points, disjoint from the blue points, be given.
If the line through any two blue points contains a red point, then the red points all lie outside the circle.
\end{lemma}

\begin{proof}
Denote the set of blue points by $B$ and the set of red points by $R$.
We define a bipartite graph $G$ with parts $B$ and $R$ by connecting $b\in B$ and $r\in R$ if the line through $b$ and $r$ does not contain any other blue points.

We first show that each $b\in B$ has degree at most $1$ in $G$.
Consider the $n-1$ lines $bb'$ where $b'\in B\setminus\{b\}$.
On each of these lines there is a distinct red point.
Thus, there is at most one red point $r$ such that $br$ contains no other blue point.
In other words, $b$ is connected to at most one red point in~$G$.

We now distinguish between two cases depending on the parity of $n$.

The case of odd $n$ is simpler.
We claim that in this case, each $r\in R$ has degree at least $1$ in $G$.
Note that each line through $r$ and a blue point contains either $1$ or $2$ blue points.
Since $n$ is odd, the number of these lines with one blue point is odd, so there is at least one.

It now follows that $G$ is a perfect matching (still when $n$ is odd).
Label the blue points $b_1,b_2,\dots,b_n$ in the order that they appear on the circle.
We consider indices to be modulo $n$.
For each $i=1,\dots,n$, let $r_i$ be the red point matched to $b_i$ in $G$.
Fix an index $i$.
Let $r_j$ be a red point on $b_{i-1}b_{i+1}$, and suppose that $r_j$ is between $b_{i-1}$ and $b_{i+1}$.
Since all lines through $r_j$ and a blue point contain two blue points except for $b_jr_j$, there can only be two more blue points apart from $b_{i-1}, b_i, b_{i+1}$.
This contradicts the assumption $n\geq 6$.
It follows that $r_j$ must be outside the circle, and $b_j=b_i$.
This shows that each red point $r_i$ lies on the line $b_{i-1}b_{i+1}$ outside the circle.
Thus all red points are outside the circle.
This finishes the case where $n$ is odd.

We now assume that $n$ is even.
Again, since each line through an $r\in R$ and a blue point contains either $1$ or $2$ blue points, it now follows that the number of lines with one blue point is even.
Thus, each $r\in R$ has even degree in $G$.
Since we have already shown that $G$ has at most $n$ edges, it follows that the number of red points with non-zero degree in $G$ is at most $n/2$.
Thus the total number of red points with zero degree is at least $n/2$.
As in the odd case, it follows from the assumption $n\geq 6$ that any red point $r$ on a line $b_{i-1}b_{i+1}$ is outside the circle.
Since there is only one blue point on $rb_i$, the degree of $r$ in $G$ is non-zero.

Furthermore, a red point can belong to at most two lines of the form $b_{i-1}b_{i+1}$, so as we go through all $n$ values of $i$, we obtain at least $n/2$ points of non-zero degree in this way.
Therefore, there are exactly $n/2$ points of non-zero degree, with each of degree exactly $2$, and each lying on two lines of the form $b_{i-1}b_{i+1}$ outside the circle.

It also follows then that there are exactly $n/2$ points of degree $0$ in $G$.
Fix an $i$ and consider a red point $r$ on $b_ib_{i+1}$.
If $r$ is inside the circle, then as before we obtain a contradiction.
Thus all of these points are outside the circle.
If one of them has non-zero degree, it has to lie on two lines of the form $b_{i-1}b_{i+1}$ as well as one of the form $b_ib_{i+1}$, which is not possible.
Thus they all have degree zero in $G$.
Since a red point can belong to at most two lines of the form $b_ib_{i+1}$, it follows that there are at least $n/2$ red points on the lines $b_ib_{i+1}$, all of degree $0$.
Thus all degree zero points are also outside the circle.
This concludes the proof for $n$ even.
\end{proof}

\end{document}
